\documentclass[11pt,a4paper]{article}

\usepackage[utf8]{inputenc}
\usepackage[T1]{fontenc}
\usepackage[brazil]{babel}
\usepackage{lmodern}
\usepackage{microtype}
\usepackage{geometry}
\usepackage{xcolor}
\usepackage{array}
\usepackage{tabularx}
\usepackage{booktabs}
\usepackage{enumitem}
\usepackage{fancyhdr}

\geometry{top=1.7cm,bottom=1.7cm,left=1.8cm,right=1.8cm}
\setlength{\parindent}{0pt}
\setlength{\parskip}{5pt}
\renewcommand{\arraystretch}{1.2}
\setlist[itemize]{leftmargin=1.5em,itemsep=2pt,topsep=3pt}
\setlist[enumerate]{leftmargin=1.7em,itemsep=3pt,topsep=3pt}

\definecolor{azul}{HTML}{173B57}
\definecolor{azulclaro}{HTML}{EAF1F7}
\definecolor{verde}{HTML}{2B6B52}
\definecolor{cinza}{HTML}{F4F5F7}
\definecolor{texto}{HTML}{27303A}

\pagestyle{fancy}
\setlength{\headheight}{15pt}
\fancyhf{}
\lhead{\textcolor{texto}{Proposta de aplicação da pesquisa}}
\rhead{\textcolor{texto}{Versão para orientação}}
\cfoot{\textcolor{texto}{\thepage}}
\renewcommand{\headrulewidth}{0.3pt}

\newcommand{\titulo}[1]{\vspace{7pt}{\large\bfseries\color{azul}#1}\par\vspace{2pt}}
\newcommand{\caixa}[1]{\begingroup\setlength{\fboxsep}{8pt}\colorbox{azulclaro}{\parbox{\dimexpr\linewidth-2\fboxsep\relax}{#1}}\endgroup}

\begin{document}

\begin{center}
  {\LARGE\bfseries\color{azul} Proposta de aplicação da pesquisa}\par
  \vspace{3pt}
  {\large Ensino introdutório de \textit{machine learning} no Ensino Médio}\par
  \vspace{5pt}
  \small Documento de trabalho para apreciação do orientador
\end{center}

\caixa{\textbf{Finalidade desta conversa.} Apresentar uma proposta objetiva de aplicação e pedir validação antes de consolidar instrumentos, sequência e procedimentos de análise. O documento descreve um planejamento, não resultados já obtidos.}

\titulo{1. Ideia central}

Pretendo desenvolver e aplicar uma sequência didática de oito aulas, apoiada por uma plataforma interativa, para ensinar fundamentos introdutórios de \textit{machine learning} a estudantes do Ensino Médio. O foco não será ensinar programação nem treinamento técnico de algoritmos. Quero investigar como os estudantes explicam a relação entre dados, regras, parâmetros, representações e resultados ao explorar situações didáticas.

\textbf{Pergunta proposta:} que mudanças e permanências podem ser identificadas nas explicações de estudantes do Ensino Médio sobre fundamentos introdutórios de \textit{machine learning} ao longo de uma sequência didática apoiada por uma plataforma interativa?

\textbf{Objetivo geral:} desenvolver e analisar uma sequência didática apoiada por uma plataforma interativa para o ensino introdutório de fundamentos de \textit{machine learning}, como conteúdo de Computação, a estudantes do Ensino Médio.

\titulo{2. O que pretendo ensinar}

O recorte proposto reúne cinco modelos, pois eles permitem comparar três finalidades: prever um valor numérico, classificar casos e formar grupos sem rótulos prévios.

\begin{center}
\small
\begin{tabularx}{\textwidth}{>{\bfseries}p{4cm} X}
\toprule
Modelo & Ideia que o estudante deverá interpretar \\
\midrule
Regressão linear & Dados, tendência, previsão numérica e erro. \\
Regressão logística & Probabilidade, limiar e decisão entre classes. \\
$k$-vizinhos mais próximos & Proximidade, exemplos rotulados e votação. \\
Árvore de decisão & Regras, perguntas e caminhos de decisão. \\
$k$-means & Agrupamento, ausência de rótulos e centróides. \\
\bottomrule
\end{tabularx}
\end{center}

Em todas as aulas, a proposta é conduzir o ciclo: \textbf{prever, experimentar, observar e explicar}. A plataforma será um recurso de mediação. Ela não substituirá a discussão conduzida pelo professor nem será tratada como causa isolada de qualquer mudança observada.

\titulo{3. Tipo de pesquisa proposto}

A proposta é uma \textbf{pesquisa aplicada}, de abordagem \textbf{predominantemente qualitativa}, com apoio de dados quantitativos descritivos. A estratégia será uma \textbf{intervenção pedagógica de grupo único}, conduzida em uma turma de Ensino Médio na qual eu atue como professor.

O estudo permitirá descrever mudanças, permanências e dificuldades nas explicações dos participantes. Não haverá grupo de controle, alocação aleatória ou pretensão de provar que a plataforma, isoladamente, causou a aprendizagem.

\newpage

\begin{center}
  {\LARGE\bfseries\color{azul} Roteiro de aplicação}\par
  \vspace{3pt}
  {\large Oito encontros de aproximadamente 50 minutos}
\end{center}

\small
\begin{tabularx}{\textwidth}{>{\centering\arraybackslash}p{1.15cm} >{\raggedright\arraybackslash}p{3.2cm} X}
\toprule
\textbf{Aula} & \textbf{Foco} & \textbf{O que será feito} \\
\midrule
1 & Dados, modelos e tipos de tarefa & Apresentar a proposta; aplicar diagnóstico individual antes do uso orientado do site; discutir previsão, classificação e agrupamento. \\
\addlinespace
2 & Regressão linear & Explorar massa, alongamento, tendência e previsão numérica. O estudante registra individualmente uma previsão antes e uma explicação depois da exploração. \\
\addlinespace
3 & Regressão logística & Explorar probabilidade, limiar e classificação. A turma discute por que uma mudança no limiar pode alterar a decisão sem modificar a probabilidade. \\
\addlinespace
4 & $k$-vizinhos mais próximos & Explorar doçura, acidez, distância, vizinhos e votação. Haverá registro individual antes e depois da exploração. \\
\addlinespace
5 & Árvore de decisão & Explorar largura, altura, regras, caminhos e profundidade. A ênfase será explicar como uma decisão é produzida. \\
\addlinespace
6 & $k$-means & Explorar agrupamento de pontos por doçura e acidez, sem categorias prévias. Haverá registro individual antes e depois da exploração. \\
\addlinespace
7 & Comparação entre modelos & Comparar finalidades, entradas, saídas, regras e limites dos cinco modelos. Preparar a turma para uma situação nova. \\
\addlinespace
8 & Síntese individual & Aplicar instrumento final individual com núcleo comparável, desafio de transferência e questionário breve de percepção. \\
\bottomrule
\end{tabularx}
\normalsize

\titulo{Organização em sala}

Cada estudante explorará a plataforma individualmente, usando os computadores do laboratório e, conforme a disponibilidade confirmada no campo, Chromebooks para completar o acesso. As produções usadas na pesquisa serão individuais. Se houver falha técnica, a aula poderá ocorrer com projeção do site e instrumentos individuais em papel.

\caixa{\textbf{Ponto de controle.} Antes da aplicação, confirmarei turma, série, componente curricular, calendário, equipamentos disponíveis e adequação dos exemplos à realidade da escola.}

\newpage

\begin{center}
  {\LARGE\bfseries\color{azul} Evidências e análise previstas}
\end{center}

\titulo{4. O que será produzido pelos estudantes}

\begin{tabularx}{\textwidth}{>{\bfseries}p{4cm} X}
\toprule
Material & Função na pesquisa \\
\midrule
Diagnóstico inicial & Registrar o ponto de partida antes do ensino orientado. Terá questões sobre dados, finalidade da tarefa, previsão, regras e limites de decisões automatizadas. Não valerá nota. \\
\addlinespace
Registros intermediários & Nos encontros 2, 4 e 6, registrar uma previsão antes da exploração e uma explicação depois dela. Esses registros ajudam a acompanhar o percurso. \\
\addlinespace
Instrumento final & Retomar, em questões equivalentes, relações avaliadas no diagnóstico; incluir uma situação inédita para examinar transferência. \\
\addlinespace
Questionário de percepção & Registrar clareza, dificuldade, uso das representações e condições de exploração da plataforma. Não será tratado como prova de aprendizagem. \\
\addlinespace
Registro docente & Anotar tempo, infraestrutura, dificuldades recorrentes, retomadas feitas e adaptações necessárias em cada aula. \\
\bottomrule
\end{tabularx}

\titulo{5. Como pretendo analisar}

As respostas serão organizadas por códigos, sem nomes dos estudantes. A análise usará uma rubrica com três níveis, de 0 a 2, em quatro dimensões:

\begin{itemize}
  \item dados e finalidade da tarefa;
  \item alteração de dado, regra ou parâmetro e sua consequência;
  \item uso de gráficos, tabelas, regras, vizinhos ou centróides na explicação;
  \item possibilidades e limites de decisões automatizadas.
\end{itemize}

Compararei apenas as questões planejadas como equivalentes entre diagnóstico e instrumento final. O resultado será apresentado por mudanças, permanências, dificuldades e respostas em branco ou incompletas. Usarei frequências absolutas e relativas somente para descrever a turma, sempre acompanhadas da leitura qualitativa das justificativas escritas.

Quando for viável, um segundo avaliador examinará uma parcela das respostas para discutir a aplicação da rubrica. O desafio de transferência será analisado separadamente, pois não terá a mesma dificuldade ou extensão do diagnóstico.

\titulo{6. Cuidados éticos e limites}

O início da produção de dados dependerá das autorizações éticas e institucionais aplicáveis. A participação será voluntária, as respostas serão anonimizadas e estudantes que não participarem da pesquisa terão acesso à atividade pedagógica equivalente. Não serão coletados dados pessoais no site.

As conclusões ficarão limitadas à turma e às condições de aplicação. A pesquisa buscará compreender o que os estudantes explicam ao longo da intervenção, sem generalizar os resultados para todo o Ensino Médio nem atribuir efeitos exclusivos à plataforma.

\newpage

\begin{center}
  {\LARGE\bfseries\color{azul} O que preciso validar na orientação}
\end{center}

\caixa{\textbf{Minha proposta para decisão:} após a conversa, ajustarei o planejamento aprovado e somente então consolidarei os instrumentos, a sequência e os procedimentos de aplicação.}

\begin{enumerate}[label=\textbf{\arabic*.}]
  \item A pergunta de pesquisa e o objetivo geral estão adequados ao foco da dissertação?
  \item A caracterização como intervenção pedagógica de grupo único, predominantemente qualitativa, é apropriada para o campo previsto?
  \item O recorte com cinco modelos é viável e pedagogicamente justificável para oito aulas?
  \item O diagnóstico, os registros intermediários e o instrumento final são suficientes para examinar mudanças nas explicações dos estudantes?
  \item A rubrica proposta e o uso de frequências descritivas são adequados para organizar os resultados?
  \item Há ajuste necessário no público, na disciplina, na duração, nos conteúdos matemáticos ou nas exigências éticas antes de seguir para o campo?
\end{enumerate}

\vfill
\begin{center}
  \textcolor{verde}{\large\bfseries Resultado esperado da reunião: uma proposta delimitada e validada para começar os ajustes finais.}
\end{center}

\end{document}
